% ============================================
% DOC-06
% Gerado de molambudos.md — build_miolo.py
% ============================================
\section*{\textbf{DOC-06} --- Relatório de Sessão III}
\markboth{DOC-06 --- Relatório de Sessão III}{}

\noindent \textit{Dr. Heitor Oliveira --- 2 de abril de 1979}

\noindent \textbf{Dr. Heitor Oliveira --- 2 de abril de 1979}
\textit{Último relatório. Sessão interrompida.}

\vspace{0.5em}\noindent\rule{\textwidth}{0.5pt}\vspace{0.5em}

\noindent \textbf{Paciente:} Joaquim Antônio Correia --- Paciente 1.260
\textbf{Tipo:} Sessão de avaliação psiquiátrica
\textbf{Duração:} 20 minutos (sessão interrompida)

\vspace{0.5em}\noindent\rule{\textwidth}{0.5pt}\vspace{0.5em}

\noindent \textbf{OBSERVAÇÃO INICIAL:}

\noindent Esta sessão foi interrompida por um incidente. O paciente apresentou alteração súbita de comportamento, e a sessão foi encerrada por segurança.

\noindent Relato a seguir:

\vspace{0.5em}\noindent\rule{\textwidth}{0.5pt}\vspace{0.5em}

\noindent \textbf{Relato:}

\noindent A sessão começou normalmente. Joaquim estava calmo. Entregou-me o diário --- algo que vinha recusando fazer havia duas semanas. Ele o colocou sobre a mesa e disse: \textit{\textquotedblleft{}Leia.\textquotedblright{}}

\noindent Perguntei se tinha certeza. Ele disse que sim.

\noindent Abri o diário. As primeiras páginas estavam escritas em caligrafia infantil, com erros de ortografia e frases curtas. As páginas descreviam a rotina no Colônia. Nada extraordinário.

\noindent Mas as últimas páginas ---

\noindent As últimas páginas estavam escritas em letra diferente. Uma letra adulta, firme, com uma caligrafia quase cursiva. Elas descreviam a \textquotedblleft{}criatura\textquotedblright{}, a \textquotedblleft{}vala\textquotedblright{}, o \textquotedblleft{}ciclo\textquotedblright{\textit{. Descreviam }este consultório\textit{. Descreviam }esta sessão}.

\noindent Havia uma passagem:

\noindent \textit{\textquotedblleft{}O médico vai ler este diário. Ele vai sentir o cheiro. Ele vai ter medo. Ele vai negar. Mas o ciclo vai começar.\textquotedblright{}}

\noindent Quando levantei os olhos, Joaquim não estava mais na cadeira.

\noindent Ele estava no canto da sala, de cócoras, com os braços enrolados em volta das pernas. Os olhos --- o olho direito --- brilhava.

\noindent Chamei-o pelo nome.

\noindent A voz que respondeu não era a dele.

\vspace{0.5em}\noindent\rule{\textwidth}{0.5pt}\vspace{0.5em}

\noindent Esta é a parte difícil de relatar.

\noindent A voz era mais grave. Mais lenta. Como se viesse de um corpo maior. Ela disse:

\noindent \textit{\textquotedblleft{}Ele não está mais aqui.\textquotedblright{}}

\noindent Perguntei quem estava falando.

\noindent \textit{\textquotedblleft{}Você sabe quem.\textquotedblright{}}

\noindent A criatura. O paciente estava em crise psicótica aguda. A entidade delirante estava se manifestando.

\noindent Tentei manter a calma.

\noindent \textit{\textquotedblleft{}O que você quer?\textquotedblright{}}

\noindent \textit{\textquotedblleft{}Nada. Só quero que você veja.\textquotedblright{}}

\noindent \textit{\textquotedblleft{}Ver o quê?\textquotedblright{}}

\noindent \textit{\textquotedblleft{}O que todo médico recusa ver.\textquotedblright{}}

\noindent Neste momento, o paciente começou a tremer. O corpo inteiro sacudiu. A boca se abriu. Os olhos reviraram. Achei que fosse uma convulsão.

\noindent Não era.

\noindent Ele parou de tremer. Ergueu a cabeça. O olho direito estava \textit{aceso} --- não é metáfora. Havia luminescência.

\noindent Ele sorriu.

\noindent \textit{\textquotedblleft{}Acha que é delírio?\textquotedblright{}}

\noindent Não respondi.

\noindent \textit{\textquotedblleft{}Acha que é trauma? Dissociação? Projeção?\textquotedblright{}}

\noindent Ainda não respondi.

\noindent \textit{\textquotedblleft{}Então me explica uma coisa, doutor.\textquotedblright{}}

\noindent Ele se levantou. Veio em minha direção. Não tocou em mim. Parou a alguns centímetros.

\noindent \textit{\textquotedblleft{}Por que você sonhou com a vala antes de ler o diário?\textquotedblright{}}

\noindent Eu não tinha contado a ele. Não tinha contado a ninguém.

\noindent \textit{\textquotedblleft{}Por que o olho dele brilha no escuro?\textquotedblright{}}

\noindent \textit{\textquotedblleft{}Isso é...\textquotedblright{}}

\noindent \textit{\textquotedblleft{}O quê? Fisiologia? Patologia?\textquotedblright{}}

\noindent Ele riu. A risada não era humana.

\noindent \textit{\textquotedblleft{}Vocês têm uma palavra para tudo. Esquizofrenia. Trauma. Delírio. Dissociação. Vocês nomeiam o que não entendem e acham que entenderam.\textquotedblright{}}

\noindent Ele se inclinou. O cheiro --- o cheiro adocicado --- era forte.

\noindent \textit{\textquotedblleft{}O nome não é a coisa, doutor. Uma coisa é não saber o nome. Outra coisa é achar que o nome é a explicação.\textquotedblright{}}

\noindent Neste momento, apaguei a luz da sala e disse que a sessão estava encerrada. O paciente não reagiu. Ele voltou para a cadeira e sentou como se nada tivesse acontecido.

\noindent Quando perguntei se ele se lembrava do que tinha acontecido, ele disse:

\noindent \textit{\textquotedblleft{}Lembro. Mas não fui eu.\textquotedblright{}}

\vspace{0.5em}\noindent\rule{\textwidth}{0.5pt}\vspace{0.5em}

\noindent \textbf{Plano Terapêutico:}

\noindent - Paciente apresenta exacerbação psicótica. Aumentar Clorpromazina para 200mg/dia.
- Solicitar avaliação neurológica (descartar epilepsia do lobo temporal).
- Sessões suspensas temporariamente.

\noindent \textbf{Nota Pessoal:}

\noindent Não vou dormir esta noite.

\vspace{0.5em}\noindent\rule{\textwidth}{0.5pt}\vspace{0.5em}

\noindent \textit{Dr. Heitor Oliveira}
\textit{Psiquiatra --- Ala de Segurança Máxima}
\textit{Hospital Colônia de Barbacena}

\noindent \textit{P.S. --- Esta noite, quando cheguei em casa, o diário estava na minha mochila. Eu não coloquei ele lá. Não sei como foi parar ali.}

\vspace{0.5em}\noindent\rule{\textwidth}{0.5pt}\vspace{0.5em}

\subsection*{ADENDO --- Redigido em 3 de abril de 1979 (1 dia após a sessão)}

\noindent Não dormi.

\noindent Tentei. Deitei. Fechei os olhos. Mas o cheiro estava no travesseiro. O cheiro adocicado. Levantei e lavei o rosto. O cheiro não saiu. Estava na minha pele.

\noindent Abri o diário de novo. Li a passagem sobre a vala. Quando levantei os olhos, eram 3:14 da manhã. O mesmo horário que Joaquim diz acordar.

\noindent Olhei para o canto do quarto. Não havia nada. Mas eu \textit{quase} vi. Uma sombra. Um movimento. Um brilho.

\noindent Fechei o diário. Coloquei dentro da gaveta. Tranquei a gaveta.

\noindent Deitei de novo.

\noindent O cheiro continuava lá.

\vspace{0.5em}\noindent\rule{\textwidth}{0.5pt}\vspace{0.5em}

\subsection*{ADENDO 2 --- 5 de abril de 1979 (72h após a sessão)}

\noindent Sonhei com a vala.

\noindent Não foi um sonho. Foi uma \textit{visita}. Eu estava deitado na pilha de corpos. Os olhos deles estavam abertos. Todos me olhando. A criatura estava em cima de mim. O rosto perto do meu. O hálito doce.

\noindent Ela disse: \textit{\textquotedblleft{}Você leu. Agora você é meu.\textquotedblright{}}

\noindent Acordei gritando.

\noindent Minha esposa acendeu a luz. Perguntou o que foi. Eu disse que não foi nada. Ela disse que eu estava branco. Eu disse que estava bem.

\noindent Mentira.

\noindent Não estou bem.

\noindent O diário está na minha cabeça. As palavras não saem. Elas se repetem. O tempo todo. Mesmo de olhos abertos.

\noindent \textit{\textquotedblleft{}Você leu. Agora você é meu.\textquotedblright{}}

\noindent Preciso dormir. Preciso parar de pensar. Preciso ---

\vspace{0.5em}\noindent\rule{\textwidth}{0.5pt}\vspace{0.5em}

\subsection*{ADENDO 3 --- 8 de abril de 1979}

\noindent O paciente 1.260 não fala mais.

\noindent Eu vou à cela dele todos os dias. Ele está sentado na cama, com o diário no colo, olhando para a parede. Não me olha quando eu entro. Não responde quando eu chamo.

\noindent Mas o olho dele --- o olho amarelo --- me segue.

\noindent Não importa onde eu esteja na cela. O olho me acompanha. Como se o olho fosse um animal separado do corpo dele.

\noindent Perguntei para a enfermeira chefe: \textquotedblleft{}O paciente 1.260 está tomando a medicação?\textquotedblright{}

\noindent Ela disse que sim. Clorpromazina 200mg. Duas vezes ao dia.

\noindent Então por que o olho dele ainda brilha?

\vspace{0.5em}\noindent\rule{\textwidth}{0.5pt}\vspace{0.5em}

\subsection*{ADENDO 4 --- 10 de abril de 1979}

\noindent Sonhei de novo.

\noindent Dessa vez eu não estava na vala. Eu estava na minha sala de consultas. O paciente 1.260 estava sentado na cadeira. Mas o rosto dele --- quando ele levantou a cabeça --- o rosto era o meu.

\noindent Acordei com o meu próprio grito.

\noindent Minha esposa dorme no sofá agora. Ela disse que eu falo dormindo. Perguntei o que eu falo. Ela disse que não quer repetir.

\noindent Perguntei se eu disse \textquotedblleft{}olho amarelo\textquotedblright{}.

\noindent Ela não respondeu.

\vspace{0.5em}\noindent\rule{\textwidth}{0.5pt}\vspace{0.5em}

\subsection*{ADENDO 5 --- 12 de abril de 1979}

\noindent Eu sei que o olho não é um delírio.

\noindent O olho brilha. Eu vi. Esta noite, fui até o hospital. Entrei na ala dos silenciados. Fiquei na porta da cela 1.260. O paciente estava dormindo.

\noindent O olho direito estava aberto.

\noindent Brilhando no escuro.

\noindent Não é possível. Não é fisiologia. Não é patologia. É outra coisa.

\noindent O que é?

\noindent O que é?

\noindent O que é?

\noindent \textit{preciso saber o que é preciso saber o que é preciso saber o que é preciso saber o que é}

\vspace{0.5em}\noindent\rule{\textwidth}{0.5pt}\vspace{0.5em}

\subsection*{ADENDO 6 --- 13 de abril de 1979}

\noindent Tranquei o diário no armário do hospital.
Joguei a chave no lixo.
Voltei para casa.
O diário estava na mesa da cozinha.

\noindent Ele não sai de perto de mim.

\noindent Ele \textit{quer} ser lido.

\noindent Ele \textit{quer} ser passado adiante.

\noindent Eu não vou mais escrever neste relatório.

\noindent Eu vou escrever no diário.

\noindent Ele me chamou.

\vspace{0.5em}\noindent\rule{\textwidth}{0.5pt}\vspace{0.5em}

\noindent \textit{Nota da direção: Dr. Heitor Oliveira não compareceu ao trabalho a partir de 14 de abril de 1979.}
\textit{Tentativas de contato telefônico fracassaram.}
\textit{Uma visita ao seu domicílio em 16 de abril constatou a residência vazia.}
\textit{O diário do paciente 1.260 não foi encontrado no hospital.}
\textit{O Dr. Oliveira permanece desaparecido.}

\vspace{0.5em}
\noindent\textit{\textbf{↪ Rotas:} \rota{MEM-16} \rotasep \rota{DOC-01}}